% Copyright (c) 2021-2027 Patrick Hall

\documentclass[fleqn]{article}
\renewcommand\refname{}
\title{Responsible Machine Learning\\\Large{Assignment 5}\\\Large{10 points}}

\usepackage{graphicx}
\usepackage{fullpage}
\usepackage{pdfpages}
\usepackage{amsmath}
\usepackage{amssymb}
\usepackage{mathtools}
\usepackage{MnSymbol}
\usepackage{enumerate}
\usepackage{setspace}
\usepackage[colorlinks, breaklinks=true]{hyperref} 
\usepackage{float}
\usepackage{caption}
\usepackage{subcaption}
\usepackage{multicol}
\usepackage{color}
\usepackage{listings}
\usepackage{csvsimple}
\usepackage{algorithm}
\usepackage{algorithmic}
\usepackage{verbatim}
\usepackage{mdframed}
\usepackage{changepage}
\usepackage[top=1in, bottom=1in, left=1in, right=1in]{geometry}

\begin{document}

\maketitle

\noindent In Assignment 5, work alone or with a group to develop and explain a tabular foundation model (TFM) for the HMDA prediction task from Assignment 1. You may use AI tools, but you are responsible for validating and understanding your results.

\setlength{\parskip}{3pt}

\section{Develop a Tabular Foundation Model.}

Use Kumo, TabPFN, TabICL, or another appropriate pretrained TFM with the target and inputs from Assignment 1. Adapt the \href{https://github.com/jphall663/GWU_rml/blob/master/lecture_6.ipynb}{Lecture 6 notebook} or use software from the \href{https://github.com/jphall663/GWU_rml/blob/master/additional_material/am6.md}{additional materials}. Document the model, package version, context selection, and settings.

Partition labeled data into training and validation sets. Draw context and any explanation-reference rows from training data; use validation data to select settings or fit a combiner or calibrator. Fit preprocessing within the appropriate fitting partition, including within cross-validation folds. Keep test data out of fitting, context, and selection; score it after choices are fixed. Report partition sizes and roles.\\

\noindent Correct TFM use, including data partitioning, is worth 3 points.

\section{Explain Your Model.}

Explain the model whose predictions you submit. Construct an additive model from frozen component-wise TFM predictions, following the lecture example, or use a technically sound explanation method from the additional-materials packages. Plot global importance or shape functions and local contributions for representative validation observations. Discuss domain plausibility and limitations. If using a surrogate, assess its fidelity to the submitted model.\\

\noindent Explainability is worth 3 points.

\section{Submit Code and Results.}

Submit your notebook and test predictions to the class evaluation server at a TBD URL. Keep your work in a public GitHub repository, or a private repository shared with the instructor (\texttt{jphall663}). Include reproducible code, partition details, validation results, explanation plots, and interpretation.

Submit a CSV with one \texttt{phat} probability column in the original test-row order, using Assignment 1's \texttt{<personal\_or\_group\_identifier>\_<model\_type>.csv} naming convention. The server ranks predictions using the \href{https://github.com/jphall663/GWU_rml/blob/master/assignment/eval.ipynb}{model evaluation notebook}. Model ranking is worth 4 points.\\

\noindent \textbf{Your deliverables are due: TBD.}

\end{document}

